% Focus: increasing transparency and cooperation at equilibrium.
% Requires graphicx; copy the two PDFs into the paper's graphics directory.

\paragraph{Transparency promotes cooperation at equilibrium.}
For the behavioral signal family, increasing transparency raises both the
maximum cooperation attainable at Nash equilibrium and the cooperation
selected by replicator dynamics. Holding the decision threshold fixed at each
of its six tested positive values, both quantities rise from zero at opacity
($t=0$) to one at full transparency ($t=1$), and neither decreases anywhere
along the sampled grid. Across the 600 adjacent transparency comparisons,
there are six increases and 594 unchanged comparisons, with no decreases.
At each grid point, exact Nash enumeration and the sampled replicator outcomes
agree: whenever a fully cooperative Nash equilibrium exists, all 1,000 sampled
initial populations converge to full cooperation. The positive relationship
therefore holds within every fixed experimental condition, as well as in the
aggregate. The observed pattern consists of thresholds and plateaus rather
than a strict increase at every transparency increment.

The aggregate in Figure~\ref{fig:behavioral-transparency-cooperation} assigns
fixed equal weights to the six tested positive decision thresholds; it is not
an average over equilibria. The within-condition results in
Figure~\ref{fig:behavioral-nash-landscape} establish that the nondecreasing
relationship does not depend on this weighting. The degenerate $\alpha=0$
condition, under which every lifted program cooperates unconditionally, is
excluded from the aggregate. Including it preserves the relationship and
changes mean cooperation at opacity from zero to $1/7$. Defecting Nash
equilibria remain available at every transparency level when $\alpha>0$;
the finding concerns attainable and evolutionarily selected cooperation.

These results use the eight-program library on the floor side, payoffs
$(T,R,P,S)=(3,2,0,-1)$, and $t\in\{0,0.01,\ldots,1\}$. The same 1,000
Dirichlet$(1)$ initial populations are reused in each distinct game. The full
707-point grid represents 62 distinct games and 62,000 stored trajectories.
The per-game pointwise 95\% Wilson interval for a cooperative basin observed
as 1,000/1,000 is $[0.9962,1]$; universal cooperation among sampled starts is
not a proof of global convergence. Step lines between sampled transparencies
are visual guides, not claims about exact continuous transition locations.

\begin{figure}[t]
  \centering
  \includegraphics[width=\linewidth]{01_transparency_and_cooperation.pdf}
  \caption{More transparency supports greater cooperation at equilibrium in
  the behavioral family. The dark curve shows maximum Nash cooperation,
  averaged with fixed equal weights over the six tested positive decision
  thresholds. Teal markers show mean cooperation at the equilibria selected
  by replicator dynamics, using the same weights. The two quantities coincide
  at all 101 sampled transparencies and rise from zero to one. Cooperation is
  measured as $\Pr[(C,C)]$. The average is over experimental settings, not
  equilibria; $\alpha=0$ is excluded because cooperation is unconditional
  there. Minimum Nash cooperation remains zero. Steps guide the eye between
  grid points separated by 0.01.}
  \label{fig:behavioral-transparency-cooperation}
\end{figure}

\begin{figure}[t]
  \centering
  \includegraphics[width=\linewidth]{02_equilibrium_landscape.pdf}
  \caption{The positive transparency effect holds within each fixed
  experimental condition. In each panel, transparency varies while caution
  is fixed. The dark line gives maximum Nash cooperation and teal markers
  give replicator-selected cooperation; both rise from zero to one and never
  decrease. Purple points show the other Nash cooperation levels, including
  persistent mutual-defection equilibria. Every Nash component here has
  constant cooperation, so all possible Nash cooperation levels are represented.
  There are no decreases across the 600 adjacent within-condition comparisons.
  Lines connect sampled values; the full-transparency endpoint at $\alpha=1$
  is shown separately.}
  \label{fig:behavioral-nash-landscape}
\end{figure}
